% !TEX root = ../algorithms.tex

\section{Трудоёмкость локальных вычислений, коммуникации и синхронизации
	при вычислениях в модели барьерно-синхронного параллелизма (BSP).
	Префиксные суммы в BSP, оптимальность.}

\begin{Definition}{Машинное слово}{}
	Машинное слово (в моделях, подобных RAM) --- это целочисленный тип данных, способный принимать все значения, необходимые для работы алгоритма.
	Машинное слово состоит из некоторого числа бит.
\end{Definition}

\subsection{Модель блочно-синхронного параллелизма (\emph{Bulk Synchronous Parallel, BSP})}

Модель \textbf{BSP} задаётся параметрами \(p\), \(g\) и \(L\), где:
\begin{itemize}
	\item \(p\) --- число компьютеров (процессоров);
	\item \(g\) --- время передачи или получения одного машинного слова;
	\item \(L\) --- время синхронизации.
\end{itemize}

Обычно предполагается, что \(g \gg 1\), \(L \gg 1\).

\begin{enumerate}
	\item Модель состоит из \(p\) компьютеров (или процессоров), пронумерованных числами от \(1\) до \(p\).
	Каждый из них представляет собой полноценную RAM-модель со своей локальной памятью.
	
	\item Компьютеры объединены сетью, и каждый из них может посылать сообщения размера \(m\) машинных слов другому компьютеру за \(gm\) единиц времени.
	
	\item После этапа локальных вычислений все компьютеры
	``синхронизируются'', ожидая завершения работы всех остальных
	компьютеров. На этом этапе доставляются все отправленные сообщения.
	Каждый компьютер, получающий \(m\) машинных слов, тратит \(gm + L\)
	единиц времени: \(gm\) на получение сообщения и \(L\) на синхронизацию.
\end{enumerate}

\subsection{Супершаги}

Алгоритм в модели BSP выполняется по \emph{супершагам}. Каждый супершаг состоит из трёх последовательных этапов:
\begin{enumerate}
	\item локальные вычисления на каждом компьютере;
	\item обмен сообщениями между компьютерами;
	\item синхронизация, на которой все компьютеры дожидаются друг друга.
\end{enumerate}

Обозначим через \(t_k^{(s)}\) время локальных вычислений, выполненных \(k\)-м компьютером на \(s\)-м супершаге.
Через \(c_k^{(s)}\) обозначим количество машинных слов, которые \(k\)-й компьютер отправил или получил на этом супершаге.

Тогда время выполнения \(s\)-го супершага определяется самым медленным вычислением, самой большой коммуникационной нагрузкой и временем синхронизации. Поэтому оно равно
\[
\max_{k=1,\dots,p} t_k^{(s)}
\;+\;
\max_{k=1,\dots,p} \bigl(g\,c_k^{(s)}\bigr)
\;+\;
L.
\]

\subsection{Стоимость BSP-алгоритма}

Пусть алгоритм работает \(\hat S\) супершагов. Тогда его время работы в модели BSP удобно измерять тремя компонентами:
\begin{enumerate}
	\item \textbf{вычисления}
	\[
	\sum_{s=1}^{\hat S} \max_{k=1,\dots,p} t_k^{(s)};
	\]
	
	\item \textbf{коммуникация}
	\[
	\sum_{s=1}^{\hat S} \max_{k=1,\dots,p} c_k^{(s)};
	\]
	
	\item \textbf{синхронизация}
	\[
	\hat S\,L.
	\]
\end{enumerate}

Во второй компоненте формально должен стоять множитель \(g\), то есть \(\sum_{s=1}^{\hat S} \max_{k} \bigl(g\,c_k^{(s)}\bigr)\), однако в асимптотических оценках его часто опускают, поскольку \(g\) считается константой и поглощается \(O\)-большим.

\paragraph{Соглашение о вводе и выводе.}
Обычно предполагается, что входные данные не лежат целиком в памяти одного компьютера. Они считаются доступными через коммуникацию с некоторым внешним хранилищем: чтобы прочитать ввод, компьютеры получают сообщения от внешнего хранилища, а чтобы выдать результат, передают его сообщениями обратно во внешнее хранилище.

\begin{Example}{Вычисление префиксных сумм в модели BSP}{}
	
	Рассмотрим задачу вычисления префиксных сумм в модели BSP. Дана последовательность из \(n\) чисел \(x_1, x_2, \dots, x_n\), где \(n \gg p\). Требуется вычислить
	\[
	y_1, y_2, \dots, y_n, \qquad y_k = \sum_{i=1}^k x_i.
	\]
	
	Для простоты будем считать, что \(p\) делит \(n\). Тогда \(k\)-й компьютер обрабатывает блок \(x_{(k-1)n/p+1}, \dots, x_{kn/p}\).
	
	\paragraph{Первый супершаг.}
	\(k\)-й компьютер считывает из внешнего хранилища свой блок длины \(n/p\).
	Вычисления не производятся, коммуникации: \(O\!\left(\frac{n}{p}\right)\).
	
	\paragraph{Второй супершаг.}
	\(k\)-й компьютер вычисляет сумму элементов своего блока и отправляет эту сумму компьютеру \(1\).
	Вычисления: \(O\!\left(\frac{n}{p}\right)\), коммуникации: \(O(p)\).
	
	\paragraph{Третий супершаг.}
	Компьютер \(1\) вычисляет префиксные суммы сумм блоков, после чего отправляет \(k\)-му компьютеру сумму всех предыдущих блоков, которую нужно прибавить к локальным префиксным суммам внутри \(k\)-го блока.
	Вычисления: \(O(p)\), коммуникации: \(O(p)\).
	
	\paragraph{Четвёртый супершаг.}
	Каждый компьютер вычисляет префиксные суммы внутри своего блока, прибавляет к ним полученную добавку и записывает результат в выходной массив.
	Вычисления: \(O\!\left(\frac{n}{p}\right)\), коммуникации: \(O\!\left(\frac{n}{p}\right)\).
	
	\begin{figure}[H]
		\centering
		\begin{tikzpicture}[
			>=Stealth, font=\small,
			box/.style={draw=black, thick, rounded corners=2pt, fill=gray!8,
				minimum width=1.7cm, minimum height=0.85cm, align=center, inner sep=2pt},
			wbox/.style={draw=black, thick, rounded corners=2pt, fill=gray!18,
				minimum width=2.6cm, minimum height=0.85cm, align=center, inner sep=2pt},
			flow/.style={->, line width=0.8pt},
			gth/.style={->, line width=0.9pt},
			sct/.style={->, line width=0.9pt, gray!55!black},
			stepp/.style={font=\small\itshape, gray!55!black, anchor=east},
			]
			
			%% columns (computers) and rows (supersteps)
			\def\cA{0} \def\cB{2.6} \def\cC{5.2} \def\cD{7.8}
			\def\yA{3.0} \def\yB{1.5} \def\yC{0.0} \def\yD{-1.6}
			
			%% step labels (left)
			\node[stepp] at (-1.8,\yA) {S1};
			\node[stepp] at (-1.8,\yB) {S2};
			\node[stepp] at (-1.8,\yC) {S3};
			\node[stepp] at (-1.8,\yD) {S4};
			
			%% column headers
			\foreach \x/\name in {\cA/$P_1$, \cB/$P_2$, \cC/$P_3$, \cD/$P_4$}
			{ \node[font=\small\bfseries] at (\x,\yA+0.85) {\name}; }
			
			%% S1: read block
			\node[box] (a1) at (\cA,\yA) {блок $x$};
			\node[box] (a2) at (\cB,\yA) {блок $x$};
			\node[box] (a3) at (\cC,\yA) {блок $x$};
			\node[box] (a4) at (\cD,\yA) {блок $x$};
			
			%% S2: local block sums T_k
			\node[box] (b1) at (\cA,\yB) {$T_1$};
			\node[box] (b2) at (\cB,\yB) {$T_2$};
			\node[box] (b3) at (\cC,\yB) {$T_3$};
			\node[box] (b4) at (\cD,\yB) {$T_4$};
			
			%% S3: prefix of block sums on P1
			\node[wbox] (pp) at (\cA,\yC) {префиксы блоков\\$P_1,\dots,P_p$};
			
			%% S4: local scan + offset -> output
			\node[box] (d1) at (\cA,\yD) {scan${}+P_1$};
			\node[box] (d2) at (\cB,\yD) {scan${}+P_2$};
			\node[box] (d3) at (\cC,\yD) {scan${}+P_3$};
			\node[box] (d4) at (\cD,\yD) {scan${}+P_4$};
			
			%% S1 -> S2 (each computer sums its own block)
			\foreach \i in {1,2,3,4} { \draw[flow] (a\i) -- (b\i); }
			
			%% S2 -> S3: gather all T_k onto P1
			\draw[gth] (b1) -- (pp);
			\draw[gth] (b2) to[out=-90,in=25] (pp.north east);
			\draw[gth] (b3) to[out=-90,in=15] (pp.north east);
			\draw[gth] (b4) to[out=-90,in=8]  (pp.north east);
			
			%% S3 -> S4: scatter offsets back
			\draw[sct] (pp) -- (d1);
			\draw[sct] (pp.south east) to[out=-25,in=90] (d2);
			\draw[sct] (pp.south east) to[out=-18,in=90] (d3);
			\draw[sct] (pp.south east) to[out=-12,in=90] (d4);
			
			%% output arrows
			\foreach \i in {1,2,3,4} { \draw[flow] (d\i) -- ++(0,-0.7); }
			\node[stepp, anchor=center, font=\footnotesize\itshape] at (\cD+1.6,\yD-0.7) {выход};
			
		\end{tikzpicture}
		
	\end{figure}
	\textit{Рис.} Поток данных в BSP-алгоритме префиксных сумм ($p=4$). \textbf{S1}: каждый $P_k$ читает свой блок длины $n/p$. \textbf{S2}: $P_k$ считает сумму блока $T_k$ и пересылает её на $P_1$ (сбор, \emph{gather}). \textbf{S3}: $P_1$ последовательно считает префиксы блок-сумм $P_k$ и рассылает $P_k$ обратно (рассылка, \emph{scatter}); именно здесь возникает узкое место $O(p)$. \textbf{S4}: каждый $P_k$ делает локальный scan и прибавляет смещение $P_k$. Число супершагов константно ($\hat S = 4$).
	\paragraph{Общее время работы.}
	Суммируя по четырём супершагам:
	\begin{itemize}
		\item вычисления: \(O\!\left(p+\frac{n}{p}\right)\);
		\item коммуникации: \(O\!\left(p+\frac{n}{p}\right)\);
		\item синхронизация: \(O(1)\), так как число супершагов константно.
	\end{itemize}
	
	Если \(p^2 \ll n\), то \(p+\frac{n}{p}=O\!\left(\frac{n}{p}\right)\). Следовательно, и вычисления, и коммуникации имеют порядок \(O\!\left(\frac{n}{p}\right)\).
	
	\paragraph{Оптимальность.}
	Алгоритм асимптотически оптимален: каждое входное машинное слово нужно хотя бы один раз прочитать, поэтому требуется не менее \(\Omega\!\left(\frac{n}{p}\right)\) коммуникаций и не менее \(\Omega\!\left(\frac{n}{p}\right)\) локальных вычислений.
\end{Example}